\documentclass[11pt,a4paper]{article}
\usepackage[margin=24mm]{geometry}
\usepackage{amsmath,amssymb,amsthm,booktabs,longtable,hyperref}
\setlength{\emergencystretch}{2em}
\newtheorem{theorem}{Theorem}\newtheorem{lemma}[theorem]{Lemma}
\newcommand{\GL}{\operatorname{GL}}\newcommand{\SL}{\operatorname{SL}}
\newcommand{\PGL}{\operatorname{PGL}}\newcommand{\PSL}{\operatorname{PSL}}
\newcommand{\gen}[1]{\langle#1\rangle}\newcommand{\diag}{\operatorname{diag}}
\title{Integral linear groups generated by three involutions}
\author{Human Compression project: N012 integral chapter}
\date{}
\begin{document}\maketitle
We recover Theorem 2.2.1 of Markovskaya's thesis~\cite{M}, including
both its linear and projective statements. The constructive matrices
are those of the source. We give explicit integral words in every
dimension, including the sign changes in its uniform construction.
The argument in dimension seven uses an additional cancellation that
is particular to the integral parameter. The finite obstructions in
dimensions three and four remain attributed external inputs.
This component does not settle the other chapters of the thesis.

Call a group \emph{admissible} if it is generated by three involutions,
two of which commute; repetitions and identity images in quotients are
allowed as in the source. In any nontrivial quotient an identity image
can instead be replaced by a suitable repeated involution. We use
column vectors and fix the conventions
\[
 t_{ij}(a)=I+aE_{ij},\qquad X^Y=YXY^{-1},\qquad
 [X,Y]=X^{-1}Y^{-1}XY.                                      \tag{1}
\]
All word identities below use (1). In particular no convention about
commutators is left implicit.

\section{Elementary generation and the excluded dimensions}
\begin{lemma}\label{roots}
The matrices $t_{ij}(1)$, $i\ne j$, generate $\SL_n(\mathbb Z)$.
Let $E$ be an integral signed permutation matrix whose permutation is
an $n$-cycle. A single unit transvection joining consecutive points
of that cycle, together with $E$, generates a group containing
$\SL_n(\mathbb Z)$. If this group contains a matrix of determinant
$-1$, it is $\GL_n(\mathbb Z)$.
\end{lemma}
\begin{proof}
Powers of $t_{ij}(1)$ give $t_{ij}(a)$ for every integer $a$.
The Euclidean algorithm, applied by elementary row additions, reduces
a primitive first column to a signed coordinate vector. The matrices
$t_{ij}(1)t_{ji}(-1)t_{ij}(1)$ interchange two coordinates with one
sign change and belong to the generated group. They allow the pivot
to be placed first, and their squares allow two signs to be changed.
Eliminating the rest of the first row and proceeding inductively
reduces a determinant-one integral matrix to the identity. This proves
the first assertion, including the final two-dimensional step.

Conjugating the given transvection by powers of $E$ gives all the
directed cycle edges, with parameters $1$ or $-1$. Inverting removes
a negative sign. For distinct $i,j,k$, direct multiplication gives
$[t_{ij}(a),t_{jk}(b)]=t_{ik}(ab)$. Following the directed cycle from
$j$ to $i$ without repeating a vertex and applying this identity
successively gives every $t_{ij}(1)$. Finally the determinant has
kernel $\SL_n(\mathbb Z)$ and image $\{1,-1\}$.
\end{proof}

We retain the following two finite exclusions from Nuzhin's
classification~\cite{N90}: $\PSL_3(2)$ and $\PSL_4(2)$ are not
admissible. These are the same precisely restricted inputs used in
the reviewed finite-linear application N011; its integral existence
input is not used here. Reduction modulo two maps each of
$\GL_n(\mathbb Z)$ and $\PGL_n(\mathbb Z)$ onto
$\GL_n(2)=\PSL_n(2)$: elementary matrices lift, and the scalar matrices
$\pm I$ reduce to $I$. Since admissibility passes to quotients, both
integral groups are excluded for $n=3,4$.

The two-dimensional boundary can be proved directly. Every nonscalar
involution in $\GL_2(\mathbb Z)$ has eigenvalues $1,-1$ over
$\mathbb Q$. Two commuting such involutions are simultaneously
diagonalizable and are equal up to sign. Consequently the projective
images of any commuting pair generate a group of order at most two.
If an admissible triple generated $\GL_2(\mathbb Z)$, its projective
quotient would therefore be generated by at most two involutions.
Reduction modulo three maps this quotient onto
$\PGL_2(3)\cong S_4$: elementary matrices and a determinant-minus-one
matrix lift, and the faithful action on the four projective points
has order $24$. But two involutions in $S_4$ generate a dihedral group
of order at most eight, since their product has order at most four.
This contradiction excludes $\GL_2(\mathbb Z)$.

For the projective positive result take
\[
 A=\begin{pmatrix}0&1\\1&0\end{pmatrix},\quad
 B=\begin{pmatrix}1&0\\0&-1\end{pmatrix},\quad
 C=\begin{pmatrix}1&1\\0&-1\end{pmatrix}.                 \tag{2}
\]
They are involutions, $AB=-BA$, and $BC=t_{12}(1)$.
Conjugating this last matrix by $A$ gives $t_{21}(1)$, so the matrices
generate $\GL_2(\mathbb Z)$; their projective images form an
admissible triple. The two boundary conclusions were previously
proved by Nuzhin--Timofeenko~\cite{NT}; the preceding arguments supply
complete proofs here rather than retaining those results as premises.

\section{The five-dimensional marking}
Let $\tau e_i=e_{6-i}$ and $\mu e_i=e_{i+1}$, with cyclic indices.
Use the source matrices
\[
 A=\begin{pmatrix}-1&0&0&0&0\\1&1&0&0&0\\0&0&1&0&0\\
 0&0&0&1&1\\0&0&0&0&-1\end{pmatrix},\qquad
 B=-\tau,\qquad C=\tau\mu,\qquad E=BC=-\mu.
                                                               \tag{3}
\]
Here $A^2=B^2=C^2=I$, $AB=BA$, and $\det E=-1$.
The following table defines each word and gives its complete value.
Factors displayed within an entry commute.
\begin{center}\small
\begin{tabular}{c|l}
word&value\\\hline
$X=[A,A^E]$&$t_{31}(-1)t_{41}(1)$\\
$Y=X^E$&$t_{42}(-1)t_{52}(1)$\\
$Z=[A,Y]$&$t_{42}(-1)t_{51}(1)t_{52}(2)$\\
$W=Z^{E^{-1}}$&$t_{31}(-1)t_{41}(2)t_{45}(1)$\\
$T=[Y,W]$&$t_{42}(-1)$\\
$V=T^{E^3}$&$t_{25}(-1)$\\
$[T,V]$&$t_{45}(1)$
\end{tabular}
\end{center}
These are integer identities, obtained from
$E_{ij}E_{kl}=\delta_{jk}E_{il}$. The last transvection joins
consecutive points of the cycle, so Lemma~\ref{roots} proves
$\gen{A,B,C}=\GL_5(\mathbb Z)$.

\section{The six-dimensional marking}
Write $p_{ij}$ for the permutation matrix interchanging $i,j$, and
$d_J$ for the diagonal matrix negative precisely on the set $J$.
For example $d_{16}$ is negative at coordinates $1,6$.
Let
\[
 A=\begin{pmatrix}-1&0&0&0&0&0\\1&1&0&0&0&0\\0&0&1&0&0&0\\
 0&0&0&0&1&0\\0&0&0&1&0&0\\0&0&0&0&0&-1\end{pmatrix},\quad
 B=p_{12}p_{34}p_{56},\quad C=p_{16}p_{25}p_{34}.        \tag{4}
\]
All three have square $I$ and determinant $-1$, and $BC=CB$.
We supply the entire word chain, since several expressions on source
p.27 have undefined subscripts or inconsistent signed-transposition
labels. The following definitions avoid those ambiguities.
\begin{center}\small
\begin{tabular}{c|l}
word&value\\\hline
$U=(A^C A)^2$&$t_{21}(1)t_{31}(1)t_{46}(-1)t_{56}(-1)$\\
$E=(A A^B)^2$&$t_{21}(-1)w_{12}d_{3456}$\\
$L=(A^B)^E$&$d_{15}t_{21}(1)p_{36}$\\
$D=(AL)^2$&$d_{3456}$\\
$D'=D^C$&$d_{1234}$\\
$F=A A^{D'}$&$d_{45}$\\
$F'=F^C$&$d_{23}$\\
$J=D'FF'$&$d_{15}$\\
$V=JL$&$t_{21}(1)p_{36}$\\
$W=V^B$&$t_{12}(1)p_{45}$\\
$K=J^C$&$d_{26}$\\
$H=(KAE)^2$&$d_{12}$\\
$I_0=HK$&$d_{16}$\\
$M=I_0 A$&$t_{21}(1)p_{45}$\\
$N=W^C$&$t_{65}(1)p_{23}$\\
$P=(MN)^2$&$t_{21}(1)t_{31}(1)t_{64}(1)t_{65}(1)$\\
$Q=(U^B)^E$&$t_{21}(-1)t_{35}(-1)t_{41}(-1)t_{65}(-1)$\\
$R=[Q,P]$&$t_{61}(1)$
\end{tabular}
\end{center}
Here $w_{12}=t_{12}(1)t_{21}(-1)t_{12}(1)$, whose $1,2$ block is
$\left(\begin{smallmatrix}0&1\\-1&0\end{smallmatrix}\right)$.
As in the previous table, multiplication of matrix units verifies
every entry exactly. Continuing gives
\[
 R^C=t_{16}(1),\qquad [A,R^C]=t_{26}(1),\qquad
 [A,R^C]^{BC}=t_{62}(1).                              \tag{5}
\]
Thus $w_{16}$ and $w_{26}$ belong to the group, where
$w_{ij}=t_{ij}(1)t_{ji}(-1)t_{ij}(1)$.
Conjugating $t_{16}(1)$ by $w_{26}$ gives $t_{12}(1)$, and conjugating
by the corresponding signed transposition on $1,2$ supplies
$t_{21}(\pm1)$. Inverting if necessary, the group contains $t_{21}(1)$.
The already available $V$ then gives
$p_{36}=t_{21}(-1)V$.

The monomial matrices
\[
             w_{16},\quad w_{26},\quad p_{36},\quad
             w_{16}^{B},\quad p_{36}^{B}                \tag{6}
\]
have permutation parts $(16),(26),(36),(25),(45)$.
These are the edges of a connected tree on the six coordinates,
so their permutation images generate $S_6$. Conjugating $t_{16}(1)$
by this monomial subgroup gives $t_{ij}(1)$ or its inverse for every
ordered pair of distinct coordinates. Lemma~\ref{roots}, together
with $\det A=-1$, now proves $\gen{A,B,C}=\GL_6(\mathbb Z)$.
In particular reduction of this exact marking gives all of
$\GL_6(3)$; this consequence requires no finite-field parameter choice.

\section{One construction for every dimension at least seven}
Put $\tau e_i=e_{n+1-i}$ and $\mu e_i=e_{i+1}$ cyclically.
Let $A$ have negative diagonal entries precisely at
$2,3,n-2,n-1$, positive entries elsewhere, and
\[
                         A_{21}=A_{n-1,n}=1.             \tag{7}
\]
All other off-diagonal entries vanish. This is the source matrix,
with its monomial generators given by
\[
\begin{array}{c|cc}
n\bmod4&B&C\\\hline
0&\tau&\tau\mu\\
1&\tau&\diag(1,\ldots,1,-1)\tau\mu\\
2&-\tau&\tau\mu\\
3&-\tau&\diag(-1,\ldots,-1,1)\tau\mu.
\end{array}                                               \tag{8}
\]
In every row $A^2=B^2=C^2=I$, $AB=BA$, and $\det(BC)=-1$.
Indeed $A$ is unchanged by reversal, while $\tau\mu$ reverses the
first $n-1$ coordinates and fixes the last, so the displayed diagonal
changes preserve its square. Counting signs gives
$(\det B,\det C)=(1,-1)$ except in the third row, where it is $(-1,1)$.

Set $E=BC$ and $\sigma=(-1)^n$. Up to a scalar sign, which has no
effect on conjugation, $E$ sends $e_i$ to $e_{i+1}$ for $i<n$ and
$e_n$ to $\sigma e_1$. Consequently
\[
 t_{ij}(a)^E=t_{i+1,j+1}(a c_i/c_j),\qquad
 c_i=1\ (i<n),\quad c_n=\sigma,                         \tag{9}
\]
with cyclic subscripts. Define
\[
 X=[A,A^E],\quad T=((XA^{E^2})^2)^E,\quad Z=[T,X],\quad
 R=[A,Z],\quad V=R^B,\quad S=[Z,V].                    \tag{10}
\]
For $n\ge8$ their exact values are
\begin{center}\small
\begin{tabular}{c|l}
$X$&$t_{31}(-1)t_{n-1,1}(\sigma)$\\
$T$&$t_{43}(-1)t_{52}(1)t_{53}(-1)t_{n3}(\sigma)$\\
$Z$&$t_{41}(1)t_{51}(1)t_{n1}(-\sigma)$\\
$R$&$t_{n-1,1}(\sigma)$\\
$V$&$t_{2n}(\sigma)$\\
$S$&$t_{21}(1)$.
\end{tabular}
\end{center}
Here is a uniform verification rather than a test of selected dimensions.
The negative diagonal positions of $A^E$ are $3,4,n-1,n$, and
its off-diagonal entries are $(3,2):1$, $(n,1):\sigma$.
For $A^{E^2}$ these positions are $4,5,n,1$, and its off-diagonal
entries are $(4,3):1$, $(1,2):1$. In the latter entry the two wrap
signs cancel. All their nonidentity
entries, and all subsequent products before the indicated conjugations,
use only the labels
\[
                       1,2,3,4,5,n-2,n-1,n.             \tag{11}
\]
These eight labels are distinct for $n\ge8$. Matrix multiplication
on them is independent of the number of unused coordinates, which
are fixed. Formula (9) transports the nonzero entries of
$(XA^{E^2})^2$ to the displayed $T$; reversal transports the two
indices of $R$ to those of $V$. Substitution in
$E_{ij}E_{kl}=\delta_{jk}E_{il}$ gives precisely the six rows,
using only $\sigma^2=1$. The accompanying checker performs this
calculation on eight abstract labels for both values of $\sigma$.
It does not infer the identity from a finite sample of dimensions.

At $n=7$, the labels $5$ and $n-2$ coincide. Direct multiplication
in seven coordinates gives the same first three rows and instead
\[
 R=t_{51}(2)t_{61}(-1),\quad
 V=t_{27}(-1)t_{37}(2),\quad
 S=t_{21}(1)t_{31}(-2).
\]
One more pair of short words separates the unit parameter:
\[
 Q=[S,S^E]=t_{31}(-1)t_{41}(2),\qquad
 [A,Q]=t_{31}(-2),\qquad S[A,Q]^{-1}=t_{21}(1).        \tag{12}
\]
This is an equality over the integers, not a division by two.
Thus in every dimension $n\ge7$ the group contains an adjacent unit
transvection and the signed cycle $E$ of determinant $-1$.
Lemma~\ref{roots} proves that it is $\GL_n(\mathbb Z)$.

\begin{theorem}
For $n\ge2$, $\GL_n(\mathbb Z)$ is admissible exactly when $n\ge5$.
The group $\PGL_n(\mathbb Z)$ is admissible exactly when $n=2$ or
$n\ge5$.
\end{theorem}
\begin{proof}
The excluded dimensions and the projective two-dimensional case were
proved in the first section. Constructions (3), (4), and (7)--(12)
cover all remaining linear dimensions. Their projective images give
the remaining projective cases, since admissibility passes to quotients.
\end{proof}

\paragraph{Attribution and verification.}
The theorem and all higher-dimensional generating triples belong to
Markovskaya--Nuzhin; the thesis attributes its general integral
construction to Nuzhin~\cite{N08}. Our contribution here is the complete
signed word calculation, the uniform eight-label verification, the
integral cancellation at $n=7$, and the shorter final generation step
at $n=6$. The two finite obstructions remain retained results.
The exact arithmetic program checks the displayed words and scope;
it is not a Lean formalization of the proof. This packet gives no
whole-thesis, new-conjecture or compression credit.

\begin{thebibliography}{9}
\bibitem{M} I.~A.~Markovskaya, \emph{Linear groups over finite fields
and over the ring of integers as automorphism groups of regular
polytopes of rank 3}, candidate dissertation, Siberian Federal
University, Krasnoyarsk, 2025, Chapter~2, pp.20--29.
\bibitem{N90} Ya.~N.~Nuzhin, \emph{Generating triples of involutions
of Chevalley groups over a finite field of characteristic 2},
Algebra i Logika \textbf{29}(2) (1990), 192--206.
\bibitem{N08} Ya.~N.~Nuzhin, \emph{On generation of the group
$\PSL_n(\mathbb Z)$ by three involutions, two of which commute},
Vladikavkazskii Matematicheskii Zhurnal \textbf{10}(1) (2008), 68--74.
\bibitem{NT} Ya.~N.~Nuzhin and I.~A.~Timofeenko,
\emph{Generating triples of involutions of linear groups of dimension
two over the ring of integers}, Vladikavkazskii Matematicheskii
Zhurnal \textbf{11}(4) (2009), 59--62.
\end{thebibliography}
\end{document}
